\section{The critical moment law and its first-field model}
\label{sec:model}

We now establish the missing input to the critical estimates: the exclusion
of the sequences in \Cref{def:prohibited-sequence}.  The exclusion has two parts.
A sufficiently steep part of the moment law produces one of the exponential
profiles already ruled out by \Cref{thm:exponential-obstruction}.  Otherwise,
after another change of scale, the first field has a nonzero limit.  Its
limiting impact equation forces both an upper growth estimate and large
values near every boundary location.  These two estimates are incompatible
with the fact that the quotient of its second derivative by the field is
holomorphic.

\begin{theorem}[Exclusion of the critical sequence]
\label{thm:sequence-exclusion}
There is no sequence satisfying \Cref{def:prohibited-sequence}.
\end{theorem}

We prove the theorem below.  Until that proof is completed, neither the
angle budget nor the polynomial bounds conditional on this theorem will
be used.

\subsection{The limiting moment law}
\label{model:moment-subsection}

Suppose, for contradiction, that a prohibited sequence is given.  Write its
primary level, depth, and normalization as
\[
 k=k_n,\qquad d=d_n=e^{-k_n},\qquad B=B_n=\mu(k_n)\longrightarrow\infty.
\]
Thus, for some \(\eps_n\to0\),
\begin{equation}
\label{model:primary-normality}
 \mu(l)\le B e^{\eps_n|l-k|}
 \qquad(l\ge k_0).
\end{equation}
Expectation without a subscript in a horizontal moment means normalized
Lebesgue measure on one period.  At physical depth \(dD\), set
\begin{equation}
\label{model:primary-moment}
 m_n(D)=\frac dB
 \left(\E\bigl(|f|^2+|q|\bigr)\right)^{1/2}
 =\frac{\mu(k-\log D)}{BD}.
\end{equation}
Each fixed compact subinterval of \(0<D<\infty\) is contained in the
domain of this definition for all sufficiently large \(n\).

\begin{lemma}[Compactness of the moment law]
\label{model:moment-compactness}
After passage to a subsequence, \(m_n\) converges locally uniformly on
\((0,\infty)\) to a positive function \(m\).  It satisfies
\begin{equation}
\label{model:moment-properties}
 m(1)=1,\qquad m(D)\le D^{-1},\qquad \log m
 \text{ is convex as a function of }D.
\end{equation}
\end{lemma}

\begin{proof}
For a periodic holomorphic function, its horizontal \(L^p\) norm,
\(p\ge1\), is log-convex in the height.  One way to see this is to regard
horizontal translates of the function as a Banach-space-valued analytic
function and apply the three-lines theorem to bounded linear functionals.
Equivalently, apply the scalar three-lines theorem followed by the dual
description of the \(L^p\) norm.  We use this for the \(L^2\) norm of \(f\)
and the \(L^1\) norm of \(q\).  If \(A\) and \(C\) are nonnegative
log-convex functions, then their sum is log-convex: the required
interpolation inequality follows from Holder's inequality applied to
\(A_0^{1-\theta}A_1^\theta+C_0^{1-\theta}C_1^\theta\).
Consequently \(\log m_n\) is convex.  A field which is identically zero
can simply be omitted from this argument.

The normalization gives \(m_n(1)=1\), and
\eqref{model:primary-normality} gives
\[
 m_n(D)\le D^{-1}\exp(\eps_n|\log D|).
\]
Thus \(\log m_n\) is uniformly bounded above on every compact
subinterval.  Convexity and its value zero at \(1\) give a lower bound
there as well.  For example, if \(D<1<D_+\), use the convexity inequality
at \(1\), viewed as a convex combination of \(D\) and \(D_+\), and the
upper bound at \(D_+\).  For \(D>1\), use a point \(D_-<1\).
Bounds on a slightly larger compact interval bound the slopes of a
convex function on the smaller one.  Arzela--Ascoli and a diagonal
subsequence therefore give a finite locally uniform limit for
\(\log m_n\).  Exponentiating proves positivity, and all the properties
in \eqref{model:moment-properties} pass to the limit.
\end{proof}

Put
\begin{equation}
\label{model:rho-definition}
 \rho(x)=-\log m(e^x).
\end{equation}
This use of \(\rho\) concerns only the moment law.  It is unrelated to
the curvature radius of the original billiard.  The function \(\rho\)
is increasing and locally Lipschitz, and
\[
 \rho(0)=0,\qquad \rho(x)\ge x.
\]
Indeed, the convex function \(\log m(D)\) tends to \(-\infty\) as
\(D\to\infty\), so each of its one-sided derivatives is strictly
negative.  With the left derivative used at exceptional points, let
\(p(x)=\rho'_-(x)>0\).  Monotonicity of the one-sided derivative of
\(\log m\) gives
\begin{equation}
\label{model:slope-inequality}
 p(v)\le e^{v-u}p(u)\qquad(v\ge u).
\end{equation}

\begin{lemma}[Contact dichotomy]
\label{lem:contact-dichotomy}
Fix \(c\in(1/2,1)\), and define
\begin{equation}
\label{model:contact-envelope}
 \rho_c(x)=cx+\inf_{v\ge x}\bigl(\rho(v)-cv\bigr).
\end{equation}
A point is called a contact if \(\rho(x)=\rho_c(x)\).
Exactly one of the following alternatives is needed.

\begin{enumerate}[label=\textup{(\roman*)}]
\item There are contacts tending to \(+\infty\) at which \(p\) is
unbounded.  In that case there exist \(x_j\to\infty\), \(P_j=p(x_j)\to
\infty\), and \(D_j=e^{x_j}\) such that
\begin{equation}
\label{model:exponential-moment-tangent}
 \frac{m(D_j(1-u/P_j))}{m(D_j)}\longrightarrow e^u
\end{equation}
uniformly on bounded real \(u\)-intervals.  For every fixed
\(\alpha_0\in(1/2,c)\), these choices can also satisfy
\begin{equation}
\label{model:contact-power-tail}
 \frac{m(D_j(1+v/P_j))}{m(D_j)}
 \le C(1+v)^{-\alpha_0}\qquad(v\ge1),
\end{equation}
with \(C\) independent of \(j\).

\item The slopes at all sufficiently large contacts are bounded.
There exist \(x_j\to\infty\), a locally Lipschitz function
\(\rho_*\) with \(\rho_*(0)=0\), and a finite \(\lambda\ge1\) such that
\begin{equation}
\label{model:rho-translate-limit}
 \rho(x_j+u)-\rho(x_j)\longrightarrow\rho_*(u)
\end{equation}
locally uniformly, and
\begin{equation}
\label{model:rho-two-ended-slope}
 \rho_*(u)=\lambda u+o(|u|)\qquad(u\to+\infty\text{ or }u\to-\infty).
\end{equation}
There are constants \(C,P'>0\), independent of \(j\), for which
\begin{align}
\label{model:rho-coarse-envelope}
 \rho(v)&\ge \rho(x_j)+c(v-x_j)-C&& (v\ge x_j),\\
\label{model:rho-fine-envelope}
 \rho(v)&\ge \rho(x_j)-P'(x_j-v)-C&& (v\le x_j).
\end{align}
Moreover \(\rho(x_j)=O(1+x_j)\).
\end{enumerate}
\end{lemma}

\begin{proof}
Since \(\rho(v)-cv\ge(1-c)v\), the infimum in
\eqref{model:contact-envelope} is attained.  The same inequality gives
\[
 x\le\rho_c(x)\le\rho(x).
\]
The running infimum is locally Lipschitz: when its lower endpoint is
moved, its value can increase only because a minimizer lay in the swept
interval, where the local Lipschitz bound for \(\rho-cx\) applies.
On an open interval containing no contact the minimizer remains to the
right, so \(\rho_c\) is affine of slope \(c\).  Everywhere,
\(\rho_c(v)-\rho_c(u)\ge c(v-u)\) for \(v\ge u\).

First suppose there are contacts \(z_j\to\infty\) with
\(p(z_j)\to\infty\).  By \eqref{model:slope-inequality},
\[
 p(x)\ge e^{-1}p(z_j)\qquad(z_j-1\le x\le z_j).
\]
For large \(j\), \(\rho(x)-cx\) is strictly increasing on this
interval.  Each point in it is itself a contact: values between that
point and \(z_j\) are larger, and values beyond \(z_j\) are at least
the value at \(z_j\) because \(z_j\) is a contact.

Set \(g=p^{-1/2}\) on this interval.  Inequality
\eqref{model:slope-inequality} says that \(e^{x/2}g(x)\) is
nondecreasing.  If \(\sigma_j\) is the supremum of \(g\) on the interval,
then \(\sigma_j\to0\), and the total variation of \(g\) there is
\(O(\sigma_j)\).  To choose a point where this small total variation
is locally well distributed, let \(\nu_j\) be its variation measure.
The elementary weak maximal inequality for intervals gives
\[
 \left|\left\{x:
 \sup_{h>0}\frac{\nu_j((x-h,x+h))}{2h}>
 \sqrt{\sigma_j}\right\}\right|=O(\sqrt{\sigma_j}).
\]
Here intervals are restricted to the original unit interval.  This
inequality follows by selecting disjoint bad intervals and covering
their union by their fivefold enlargements.  Choose \(x_j\) outside
this bad set in the middle third of the unit interval.  Then, whenever
\(|h|\) is smaller than its distance to the endpoints,
\[
 |g(x_j+h)-g(x_j)|\le
 2\sqrt{\sigma_j}|h|.
\]
Taking \(|h|\le g(x_j)=P_j^{-1/2}\) shows that
\[
 \frac{p(x_j+h)}{P_j}\longrightarrow1
 \quad\hbox{uniformly for }|h|\le P_j^{-1/2}.
\]
Integration, with
\(\log(1-u/P_j)=-u/P_j+O(P_j^{-2})\), proves
\eqref{model:exponential-moment-tangent}.

For the tail estimate write \(z=\log(1+v/P_j)\).  If
\(0\le z\le P_j^{-1/2}\), integration of the last slope estimate gives
\(\rho(x_j+z)-\rho(x_j)\ge P_jz/2\), which is at least a fixed
positive multiple of \(v\) in this range.  At
\(z=P_j^{-1/2}\) the increase is at least a fixed multiple of
\(\sqrt{P_j}\).  This endpoint is still a contact.  For larger \(z\),
the contact inequality adds at least
\(c(z-P_j^{-1/2})\).  Since
\[
 \log(1+v)\le \log P_j+z+O(1)
\]
and \(\alpha_0<c\), the increase is at least
\(\alpha_0\log(1+v)-O(1)\), uniformly for all \(v\ge1\).
Exponentiating proves \eqref{model:contact-power-tail}.

Now suppose that \(p\) is bounded on contacts in a tail.  Almost
everywhere on the contact set the derivatives of \(\rho\) and
\(\rho_c\) agree, since the two locally Lipschitz functions agree
there.  Off that set the derivative of \(\rho_c\) equals \(c\).
Consequently \(\rho_c\) has a uniform Lipschitz constant \(K_1\) on
some tail.  Extend it to the left with a fixed Lipschitz constant.
Consider the anchored translates
\[
 u\longmapsto\rho_c(x+u)-\rho_c(x).
\]
They lie in a compact space of anchored Lipschitz functions, with
the topology of uniform convergence on compact intervals.  Average
their point masses over \(x\in[x_0,x_0+R]\).  A weak limit as
\(R\to\infty\) is invariant under every anchored translation
\[
 (\mathcal T_h q)(u)=q(u+h)-q(h).
\]
Indeed translating the averaging interval changes its normalized
measure by at most \(2|h|/R\).

The mean increment at \(1\) of this invariant law is at least \(1\).
To check this, telescope the integral:
\[
 \frac1R\int_{x_0}^{x_0+R}
       \bigl(\rho_c(x+1)-\rho_c(x)\bigr)\,\dd x
 =
 \frac1R\left(
 \int_{x_0+R}^{x_0+R+1}\rho_c-
 \int_{x_0}^{x_0+1}\rho_c\right),
\]
and use \(\rho_c(x)\ge x\). Birkhoff's pointwise ergodic theorem
\cite[Theorem~6.2.1]{Durrett2019} for the invertible map
\(\mathcal T_1\), applied in both directions to the bounded unit
increments, now provides a
path \(q_*\) in the support of the law for which
\begin{equation}
\label{model:envelope-ergodic-path}
 \frac{q_*(u)}u\longrightarrow\lambda\ge1
 \qquad(u\to+\infty\text{ or }u\to-\infty).
\end{equation}
The two limits agree because they are the same invariant conditional
expectation of the unit increment.  Bounded increments between
successive integers extend the conclusion to real \(u\).
Such a path is approximated locally uniformly by translates with
centers \(x_j\to\infty\): a neighborhood visited only for bounded
\(x\) has zero mass in every limiting averaged measure.

We next replace \(\rho_c\) by \(\rho\).  To the right of any fixed
compact \(u\)-interval, \eqref{model:envelope-ergodic-path} gives a
fixed interval \([v,w]\) on which the increment of \(q_*\) exceeds
\(c(w-v)\).  The approximating translates of \(\rho_c\) must
therefore have a contact in \([x_j+v,x_j+w]\).  Monotonicity of
\(\rho\), together with \(\rho\ge\rho_c\), bounds
\(\rho(x_j+u)-\rho_c(x_j)\) above and below on every compact interval.
It also bounds the local slopes of \(\rho\): from
\eqref{model:slope-inequality},
\[
 p(x)\le e\int_{x-1}^x p(v)\,\dd v
       =e\bigl(\rho(x)-\rho(x-1)\bigr).
\]
Thus a subsequence of the anchored translates of \(\rho\) has a
locally uniform limit \(\rho_*\), and
\(\rho(x_j)-\rho_c(x_j)=O(1)\).

For completeness, this limit has the same two-ended slope as \(q_*\).
Fix a small \(\eta>0\).  For every sufficiently large \(|u|\),
the increment of \(q_*\) on \([u,u+\eta|u|]\) exceeds
\(c\eta|u|\), by \eqref{model:envelope-ergodic-path} and \(\lambda>c\).
The same contact argument, followed by passage to the limit, bounds
the excess of \(\rho_*(u)\) over \(q_*(u)\), up to a fixed additive
constant, by \(K_1\eta|u|\).  The lower bound
\(\rho\ge\rho_c\) bounds the deficit by a fixed constant.
Letting \(\eta\downarrow0\) proves
\eqref{model:rho-two-ended-slope}.

Finally, the lower increment bound for \(\rho_c\), its uniform
Lipschitz bound, and the bounded difference at \(x_j\) give
\eqref{model:rho-coarse-envelope} and
\eqref{model:rho-fine-envelope} as long as \(v\) remains in the tail.
For the remaining left range, use \(\rho(v)\ge v\), increasing
\(P'>\max(K_1,1)\) and the additive constant if necessary.
The same Lipschitz bound gives \(\rho(x_j)=O(1+x_j)\).
\end{proof}

\subsection{The steep-contact alternative}
\label{model:steep-subsection}

\begin{lemma}
\label{model:steep-excluded}
Alternative \textup{(i)} of \Cref{lem:contact-dichotomy} is impossible
for the moment law of a prohibited sequence.
\end{lemma}

\begin{proof}
For each fixed \(j\), take the original index \(n\) sufficiently large
and set
\begin{equation}
\label{model:contact-physical-scales}
 h_n=H-dD_j,\qquad
 \ell_n=\frac{dD_j}{P_j},\qquad
 m_0=\frac Bd\,m_n(D_j).
\end{equation}
The diagonal can be chosen so that
\[
 h_n\to H,\qquad \ell_n\to0,\qquad
 m_0\ell_n\to\infty.
\]
For the last assertion, first keep \(j\) fixed:
\(m_0\ell_n=B D_jm_n(D_j)/P_j\to\infty\).
The distance from \(h_n\) to the upper shape wall, in units of
\(\ell_n\), is \(P_j\to\infty\); the lower wall also escapes in
these units.

Both \eqref{model:exponential-moment-tangent} and
\eqref{model:contact-power-tail} can be transferred to the actual
moments on the required ranges.  Here is the tail detail.  Fix
\(\alpha_0\in(1/2,c)\), and require
\(\eps_n<(1-\alpha_0)/2\).  For each fixed \(j\), choose a finite
cutoff \(D_{\rm cut}\) large enough that the primary bound
\(m_n(D)\le D^{-1+\eps_n}\) is below a fixed multiple of the
right side of \eqref{model:contact-power-tail}, with the
normalizing factor \(m_n(D_j)\), whenever \(D\ge D_{\rm cut}\).
Such a cutoff exists because \(1-\eps_n>\alpha_0\) and
\(m_n(D_j)\to m(D_j)>0\).  On the finite interval below the cutoff
use locally uniform moment convergence. Choose the fixed physical
cutoff \(0<\delta_b<H/2\) small enough that the enlarged depth range
used for submean disks has logarithmic level at least \(k_0\).
Thus all the coarser depths up to \(\delta_b\)
are controlled, with a constant independent of \(j\).

On the row of height \(h_n\), define
\[
 H_x=\frac{|f(x+\ii h_n)|^2+|q(x+\ii h_n)|}{m_0^2}.
\]
Its horizontal mean is \(1\).  Sample \(x\) with probability density
\(H_x\) relative to normalized horizontal measure, and set
\(m_x=m_0\sqrt{H_x}\); the zero-density set \(H_x=0\) can be
discarded.  In the coordinate
\[
 S=\frac{s-x-\ii h_n}{\ell_n},
\]
let
\[
 \mathcal F_x(S)=\frac{f(s)}{m_x},\qquad
 \mathcal Q_x(S)=\frac{q(s)}{m_x^2}.
\]
There is an exact cancellation in the expectation:
\begin{equation}
\label{model:biased-cancellation}
 \E_H\bigl(|\mathcal F_x(S)|^2+|\mathcal Q_x(S)|\bigr)
 =
 \frac{\E\bigl(|f|^2+|q|\bigr)
       \text{ on the row }h_n+\ell_n\Im S}{m_0^2}.
\end{equation}
The denominator \(H_x\) cancels against the sampling density,
and the remaining integral is invariant under the horizontal
translation \(x\mapsto x+\ell_n\Re S\).  The transferred local
moment law therefore makes \eqref{model:biased-cancellation}
converge locally uniformly to \(e^{2\Im S}\).

Submean estimates applied before taking expectation give the same
bounds for suprema on smaller tiles.  In a cone of fixed aperture
\(L\), a dyadic band with depth \(D\ge1\) requires only a number of
depth-sized tiles depending on \(L\).  The expected squared size of
the fields on this band is at most \(C_LD^{-2\alpha_0}\).
Choose one \(\alpha\in(1/2,\alpha_0)\), independently of \(L\).
Summation of
\(D^{2\alpha}C_LD^{-2\alpha_0}\) over dyadic bands proves tightness
of the random constants in the pointwise bounds
\begin{equation}
\label{model:sampled-coarse-bound}
 \max(|\mathcal F_x(S)|,\sqrt{|\mathcal Q_x(S)|})
 \le C_L(-\Im S)^{-\alpha}
\end{equation}
for \(1\le-\Im S\le\delta_b/\ell_n\) and
\(|\Re S|\le L(-\Im S)\).
Also
\[
 \Prob_H(H_x\le a)=\E\bigl(H_x\ind_{\{H_x\le a\}}\bigr)\le a.
\]
Since \(m_0\ell_n\to\infty\), it follows that
\(m_x\ell_n\to\infty\) in probability.

Now consider the nonnegative subharmonic function
\[
 U_x(S)=|e^{\ii S}\mathcal F_x(S)|^2
       +|e^{2\ii S}\mathcal Q_x(S)|.
\]
Its mean tends locally uniformly to \(1\).  If a nonnegative smooth
cutoff equals \(1\) on a fixed compact set, integration against its
Laplacian shows that the expected Riesz mass of \(U_x\) on that set
tends to zero: the limiting integral of the Laplacian of the cutoff
is zero.  The mass itself is nonnegative, so it tends to zero in
probability.

We may therefore select deterministic centers and a subsequence with
all the following properties: local analytic bounds on an exhaustion
of the plane; the bounds \eqref{model:sampled-coarse-bound} for every
fixed aperture; \(m_x\ell_n\to\infty\); and vanishing Riesz mass on
every fixed compact set.  To make the quantifier explicit, assign
summable failure probabilities to countably increasing apertures
and compact sets, choose their bound constants once, and then take
the sequence index sufficiently large for the remaining probability
conditions.  The union of the bad events has probability less than
one.  This gives fixed eventual constants for each aperture, not
constants growing with the sequence index.

The resulting entire limits \((\mathcal F,\mathcal Q)\) obey
\[
 |e^{\ii S}\mathcal F(S)|^2
       +|e^{2\ii S}\mathcal Q(S)|
 \quad\hbox{harmonic}.
\]
Each summand is subharmonic, so each is harmonic.  An analytic
function with harmonic squared modulus is constant.  An analytic
function with harmonic modulus is also constant: away from its
zeros its Laplacian is \(|G'|^2/|G|\), and the assertion follows
by analytic continuation.  Thus
\[
 \mathcal F(S)=c_1e^{-\ii S},\qquad
 \mathcal Q(S)=c_2e^{-2\ii S}.
\]
At \(S=0\) the normalized sum is exactly \(1\), so the limit is
nonzero.  Moreover
\[
 \mathcal F_x''=(m_x\ell_n)^2\mathcal Q_x\mathcal F_x.
\]
Compact convergence and \(m_x\ell_n\to\infty\) imply
\(\mathcal Q\mathcal F=0\).  The limit is consequently one of
the two nonzero pure exponential profiles in
\Cref{thm:exponential-obstruction}.  All its finite-wall hypotheses
were checked above, including a single exponent \(\alpha>1/2\)
for every fixed aperture.  That theorem gives the contradiction.
\end{proof}

\subsection{Re-centering and selecting a nonzero first field}
\label{model:selection-subsection}

We are therefore in alternative \textup{(ii)} of
\Cref{lem:contact-dichotomy}.  Choose its \(x_j,\rho_*,\lambda\).

\begin{lemma}[The second scale]
\label{model:selected-fields}
There are depths \(d'\to0\), levels \(K=-\log d'\to\infty\),
normalizations \(B'=\mu(K)\to\infty\), and horizontal centers
\(\xi_n\) with the following properties.  Set
\begin{equation}
\label{model:selected-normalization}
 F_n(S)=\frac{d'}{B'}f(\xi_n+\ii H+d'S),\qquad
 Q_n(S)=\left(\frac{d'}{B'}\right)^2q(\xi_n+\ii H+d'S).
\end{equation}
Then \(F_n\to F\not\equiv0\) and \(Q_n\to0\) locally uniformly on
\(\mathbb H_-=\{S:\Im S<0\}\), and \(F''/F\) is holomorphic there.
The global normalized moment laws converge locally uniformly to
\begin{equation}
\label{model:star-moment}
 m_*(D)=\exp\bigl(-\rho_*(\log D)\bigr),\qquad
 \log m_*(D)=-\lambda\log D+o(|\log D|)
\end{equation}
at both ends of \(0<D<\infty\).  For fixed
\(1-c<\beta<1/2\) and some fixed \(a>0\),
\begin{equation}
\label{model:new-global-envelope}
 \mu(l)\le C B'
 \exp\bigl(\beta(K-l)_++a(l-K)_+\bigr)
 \qquad(l\ge k_0),
\end{equation}
with \(C\) independent of the sequence index.  Finally, one can fix
\(\alpha\in(1/2,1-\beta)\) so that, for every fixed aperture \(A\),
\begin{equation}
\label{model:selected-cones}
 \max(|F_n(S)|,\sqrt{|Q_n(S)|})
 \le C_A(-\Im S)^{-\alpha}
\end{equation}
when \(1\le-\Im S\le\delta_b/d'\) and
\(|\Re S|\le A(-\Im S)\), eventually for that aperture.
\end{lemma}

\begin{proof}
For each fixed \(j\), take the original index \(n\) sufficiently
large and put
\[
 d'=d e^{x_j},\qquad K=k-x_j,\qquad B'=\mu(K).
\]
The identity
\[
 \frac{B'}B=e^{x_j}m_n(e^{x_j})
\]
shows that \(K,B'\to\infty\) are compatible with this order of
choices. We also impose \(s_n>x_j+3\), so the ancestor cutoff lies
strictly below the new target level. The new normalized moment law is
\[
 \frac{m_n(e^{x_j}D)}{m_n(e^{x_j})},
\]
which tends, after the diagonal selection, to
\eqref{model:star-moment}.

We spell out the global envelope because its constants must remain
uniform after diagonalization.  With \(v=k-l\),
\begin{equation}
\label{model:moment-ratio-change}
 \frac{\mu(l)}{B'}=
 e^{v-x_j}\frac{m_n(e^v)}{m_n(e^{x_j})}.
\end{equation}
On a finite \(v\)-range, the limiting estimates
\eqref{model:rho-coarse-envelope} and
\eqref{model:rho-fine-envelope} imply an upper exponent \(1-c\)
for \(v\ge x_j\) and \(P'-1\) for \(v\le x_j\).
Choose \(\beta>1-c\) and \(a>\max(P'-1,0)\).
Since \(\rho(x_j)=O(1+x_j)\), primary normality bounds the logarithm
of \eqref{model:moment-ratio-change} by
\[
 C(1+x_j)+\eps_n|v|.
\]
If \(\eps_n\) is small compared with \(\beta\) and \(a\), this is
absorbed by the chosen envelope for
\(|v|\ge C_0(1+x_j)\), with \(C_0\) fixed sufficiently large.
On the remaining finite interval take \(n\) large enough for
moment convergence with a fixed multiplicative error.  This proves
\eqref{model:new-global-envelope} with one constant \(C\).

Let \(k_a=k-s_n\) and \(\omega,\nu\) be the ancestors and measures
in \Cref{def:prohibited-sequence}.  To apply
\Cref{prop:weighted-bootstrap}, take an old coarse exponent
\(0<\beta_0<\beta\).  For \(l\le k_a+3\), the ratio of the new
coarse normalization to the old one equals
\begin{equation}
\label{model:ancestor-domination}
 \frac{B'e^{\beta(K-l)}}{Be^{\beta_0(k-l)}}
 =
 \frac{B'}B e^{-\beta x_j}
 e^{(\beta-\beta_0)(k-l)}
 \ge
 \frac{B'}B e^{-\beta x_j}
 e^{(\beta-\beta_0)(s_n-3)}.
\end{equation}
For fixed \(j\), the last expression tends to infinity.  We can
therefore make it at least one before proceeding to the next \(j\).
The same comparison, squared, applies to the absolute ancestor
integral of \(P\).  Thus the coarse hypotheses of the weighted
bootstrap hold with constants independent of \(j\).  The global
envelope already has this uniformity, so the resulting weighted
\(L^2\) bounds do as well.

For each fixed \(j\), the new \(q\)-field on a fixed compact set is
the primary \(q\)-field on another fixed compact set, multiplied by
the finite factor \(m_n(e^{x_j})^{-2}\).  The primary convergence
to zero in \(\nu\)-probability can therefore be imposed to any
prescribed tolerance before advancing \(j\).  Hence the new
\(q\)-field tends to zero in probability on compact sets.

Weighted submean estimates now give tightness of the analytic fields
on every compact subset of \(\mathbb H_-\).  They also give
pointwise cone tightness.  At depth ratio \(D\ge1\), the coarse
part of \eqref{model:new-global-envelope} bounds the weighted
field norm by \(CD^{-(1-\beta)}\).  Summing depth-sized tile
suprema over dyadic bands, with any fixed
\(\alpha\in(1/2,1-\beta)\), proves tightness of the constants in
\eqref{model:selected-cones}.  Fine tiles use the bounded change of
the density of \(\nu\) on shifts of order the ancestor depth.
For tiles coarser than that depth, the necessary tile estimates
are precisely the explicit coarse hypotheses of the weighted
bootstrap.  Thus no translation invariance of \(\nu\) is required.

Choose deterministic centers on events with these compact and cone
bounds and with the \(q\)-field small, using summable exceptional
probabilities as in the preceding proof.  Extract a locally
uniform limit \((F,0)\).  If \(F\) were zero, its compact
smallness and the coarser bound with exponent \(\alpha>1/2\)
would violate \Cref{thm:no-trivial-profile}; the global
normalization at the target is \(B'=\mu(K)\), as that theorem
requires.  Therefore \(F\not\equiv0\).

Each quotient \(F_n''/F_n=(d')^2q\) is holomorphic.  Around any
zero of \(F\), choose a small circle containing no other zero and
no zero on the circle.  The quotients converge on that circle to
\(F''/F\), so maximum modulus bounds them in its interior.
Their limit removes the apparent singularity of \(F''/F\).
This proves the last asserted regularity.  In particular all zeros
of \(F\) are simple and \(F''\) vanishes at them: otherwise the
analytic equation \(F''=(F''/F)F\) with initial data
\(F=F'=0\) would force \(F\equiv0\).
\end{proof}

\subsection{The limiting impact equation}
\label{model:equation-subsection}

Use the scales and centers of \Cref{model:selected-fields}.  The phase
coordinate and parameter are now
\begin{equation}
\label{model:phase-scaling}
 Y=\frac{y-\xi_n-\ii H}{d'},\qquad
 T=\frac{B't}{d'},\qquad
 S_n(Y,T)=\frac{\mathfrak s(y,t)-\xi_n-\ii H}{d'}.
\end{equation}
The cone estimates in this subsection refer to depths
\(D=-\Im Y\) or \(D=-\Im S\), as indicated.

\begin{proposition}[The pure first-field model]
\label{prop:pure-model}
There is a limiting impact coordinate \(S(T,Y)\) with the following
properties.

\begin{enumerate}[label=\textup{(\roman*)}]
\item For every fixed aperture and every bounded complex parameter
range, it is jointly holomorphic sufficiently deep in the corresponding
cone and is the locally uniform limit of \eqref{model:phase-scaling}.
There,
\begin{equation}
\label{model:deep-anchor}
 S(T,Y)-Y=O\bigl(|T|^2D^{1-2\alpha}\bigr),
\qquad D=-\Im Y\longrightarrow\infty.
\end{equation}

\item Put
\begin{equation}
\label{model:energy}
 J=S_Y,\qquad z=\frac{TF(S)J}{2},\qquad
 g(z)=\frac{3(\sin z-z\cos z)}{z^3},\quad g(0)=1.
\end{equation}
The limiting equation is
\begin{equation}
\label{model:energy-equation}
 J^3g(z)=1.
\end{equation}

\item For every fixed \(\eps>0\), set \(r=1/(\lambda+\eps)\).
There are constants \(C_\eps,t_0>0\) such that, for every real
\(|T|<t_0\), the model is holomorphic in \(Y\) and maps
\begin{equation}
\label{model:real-domains}
 \{Y:\Im Y<-C_\eps|T|^r\}\ \longrightarrow\ \mathbb H_-.
\end{equation}
It satisfies \(S(0,Y)=Y\), obeys
\eqref{model:energy-equation} throughout this domain, and depends
continuously on real \(T\), locally uniformly on strict common
interiors.

\item At every strict interior point of \eqref{model:real-domains}
with \(T\) real, these functions have a joint holomorphic germ in
\((T,Y)\).  The germs agree when continued along vertical \(Y\)-paths
from the deep model.
\end{enumerate}
\end{proposition}

\begin{proof}
For a bounded complex \(T\)-range, apply \Cref{thm:cone} to
\eqref{model:selected-cones} at a sufficiently large fixed depth
ratio.  Its smallness condition holds because
\(|T|C_AD^{-\alpha}\) is small when \(D\) is large.  The physical
parameter is \(t=Td'/B'\); relative to a fixed positive multiple
of \(d'\), its size tends to zero.  Thus the constant bounding the
parameter-to-depth ratio in that theorem can be fixed, for example
at one.  The resulting passages are attached vertically to the
original lift and satisfy
\[
 S_n-Y=O\bigl(|T|^2D^{1-2\alpha}\bigr)
\]
on each fixed deep cone.  They stay in the shape strip.  The bound
and holomorphic compactness give joint subsequential limits.

We derive their equation before proving uniqueness.  Write
\(h=T/B'\), and use the rescaled step
\[
 J_n(Y)=\frac{S_n(Y+h,T)-S_n(Y,T)}h.
\]
On the fast segment based at a scaled shape coordinate \(S\), the
fields entering the normalized chord action are
\[
 T F_n(S+hu),\qquad T^2Q_n(S+hu).
\]
They converge analytically, including on small thickenings, to the
constant fast pair \((TF(S),0)\).  The kernel division in
\Cref{lem:kernel-division} and the frozen first-field formula
\eqref{eq:frozen-F-energy} therefore give the local action limit
\[
 p(S,J)=\frac{J}{(TF(S))^2}
       -\frac{2\sin(TF(S)J/2)}{(TF(S))^3},
\]
with its removable interpretation at \(TF(S)=0\).
This passage is initially made where the deep near-identity
branches have nonvanishing kernels.

In scaled coordinates the exact stationary equation reads
\[
 p_{n,J}(S,J)-h\,p_{n,S}(S,J)
       =p_{n,J}(S_-,J_-).
\]
The preceding and current states are evaluations of the same
holomorphic functions at phase arguments differing by \(h\).
Dividing their difference by \(h\) and using analytic convergence
gives
\[
 \frac{\partial}{\partial Y}p_J=p_S,\qquad J=S_Y.
\]
In particular no prescribed rate of convergence relative to \(h\)
is needed.  Differentiating \(E=Jp_J-p\) along the limiting
solution gives
\[
 E_Y=J\frac{\partial p_J}{\partial Y}-p_SS_Y=0.
\]
The deep anchor implies \(J\to1\), while \(F(S)\to0\), so
\(E=1/12\).  The frozen energy identity
\[
 E=\frac{2(\sin z-z\cos z)}{(TF(S))^3}
\]
is exactly \eqref{model:energy-equation}.

Near zero field, \eqref{model:energy-equation} determines an
analytic branch \(J=j(TF(S))\) with \(j(0)=1\), since its
\(J\)-derivative at zero field and \(J=1\) is \(3\).
It has \(j(w)^{-1}-1=O(w^2)\).  Define the inverse coordinate
\(\mathcal Y_T\) deep in the cone by
\[
 \mathcal Y_T'(S)=j(TF(S))^{-1},\qquad
 \mathcal Y_T(S)-S\longrightarrow0
       \quad\hbox{down the cone}.
\]
The normalized primitive exists because its derivative minus one
is \(O(|T|^2D^{-2\alpha})\) and \(2\alpha>1\).  It can be
constructed by integration from a point at depth \(R\), then
letting \(R\to\infty\); changing the horizontal coordinate of the
starting point contributes \(O(R^{1-2\alpha})\), which vanishes.
Along every limiting solution,
\[
 \frac{\partial}{\partial Y}\mathcal Y_T(S(T,Y))=1.
\]
The anchor fixes the additive constant to zero.  Hence the solution
is the unique nearby inverse of \(\mathcal Y_T\) sufficiently deep.
All deep subsequential limits agree, and the convergence and
holomorphic dependence asserted in \textup{(i)} follow.

We next obtain the domains for real \(T\).  Fix
\[
 \frac12<c_1<1-\beta
\]
and a sufficiently large fixed number \(C_0\).  For a small fixed
depth ratio \(D>0\), the global moment laws on ratios
\(v\ge D/C_0\) have the bound
\begin{equation}
\label{model:width-moment-cone}
 C_\eps D^{-\lambda-\eps}(D/v)^{c_1}.
\end{equation}
For \(D/C_0\le v\le1\), first use
\(m_*(v)\le C_\eps v^{-\lambda-\eps/2}\), then take the original
index sufficiently large at this fixed \(D\).
For \(v\ge1\), use
\eqref{model:new-global-envelope}, whose field exponent is
\(1-\beta>c_1\).  These give \eqref{model:width-moment-cone}
with constants independent of \(D\), although the required
sequence index may depend on \(D\).

Submean and dyadic tile summation, with a slightly smaller exponent
still greater than \(1/2\), provide at least one horizontal center
whose actual fields obey the corresponding pointwise cone bound.
This step uses uniform horizontal sampling and the global moments;
the good center need not be \(\xi_n\).  Apply \Cref{thm:cone}
at that center with
\[
 |T|\le c_\eps D^{\lambda+\eps}.
\]
After choosing its small constant and, if necessary, replacing the
base depth by a fixed smaller multiple, the passage reaches phase
depth \(d'D\).  The real no-patch transfer, using
\Cref{thm:no-patch} as after \Cref{thm:cone}, consequently gives
the global width bound
\begin{equation}
\label{model:real-width-bound}
 \Delta(Td'/B')\le d'D
 \quad\hbox{for }|T|\le c_\eps D^{\lambda+\eps},
\end{equation}
eventually for each fixed \(D\).  Recall that
\(\Delta(t)=H-h_X(t)\).
This is the point where a successful local passage is converted
to a full horizontal domain; local analytic continuation alone
would not give \eqref{model:real-width-bound}.

The full-width domains in \eqref{model:real-width-bound} exhaust
\eqref{model:real-domains}, with a strict margin on compact sets.
The shape constraint gives \(\Im S_n<0\), and the deep anchor
gives a fixed interior normalization.  Holomorphic maps to a
half-plane with such an anchor form a normal family, as is seen by
mapping the half-plane to the unit disk.  Every limit agrees with
the unique deep solution.  This defines \(S(T,\cdot)\) on
\eqref{model:real-domains}; the energy identity propagates there
by the identity theorem.
The same argument for a sequence of real parameters \(T_j\to T\)
proves local uniform continuity: on a common strict interior,
every convergent subsequence has the same deep germ and hence the
same limit.  This also treats \(T=0\).

It remains to justify joint germs at all real parameters.
Equation \eqref{model:energy-equation} first shows that \(J\ne0\).
At fixed \(w=TF(S)\), its \(J\)-derivative is
\begin{equation}
\label{model:twist-derivative}
 \frac{\partial}{\partial J}\bigl(J^3g(wJ/2)\bigr)
       =3J^2\frac{\sin z}{z}.
\end{equation}
It can vanish only when \(z=n\pi\), \(n\in\Z\setminus\{0\}\).
The equivalent entire equation
\begin{equation}
\label{model:cubic-energy}
 (TF(S))^3=24(\sin z-z\cos z)
\end{equation}
then has a nonzero right side.  Differentiating in \(Y\) gives
\[
 3T^3F(S)^2F'(S)J=24z\sin z\,z_Y,
\]
and hence \(F'(S)=0\) at such a point.
The function \(F\) is nonconstant, because it is nonzero and
decays down the cones.  Its critical points are therefore
countable.  At each such point and each nonzero integer \(n\),
\eqref{model:cubic-energy} allows at most three values of \(T\).
Thus the exceptional real parameters form at most a countable set.

This countable exception is removable.  Fix an exceptional real
\(T\), a strict interior point, and a nearby sequence of
nonexceptional real parameters.  On a common connected interior
\(Y\)-domain, continuity gives local uniform convergence of
\(z(Y,T_j)\) to \(z(Y,T)\).  For each fixed nonzero integer \(n\),
Hurwitz's theorem says that \(z(Y,T)-n\pi\) is either zero-free
or identically zero.  The latter alternative is impossible on
the common domain: analytic continuation to its deep part gives
\(z\to0\).  Exhausting the strict interior proves omission of all
the exceptional values for every real \(T\).

Consequently \eqref{model:twist-derivative} is nonzero along every
compact vertical path in the real-parameter domain.  Solve the
regular implicit equation as an analytic first-order ODE in \(Y\),
starting with the joint deep data.  A finite chain of local ODE
neighborhoods along the path gives a joint germ in \((T,Y)\) at
its endpoint.  The constructed germ agrees with the given
real-parameter functions by their continuity and local ODE
uniqueness.  Germs obtained on overlapping vertical paths agree
by the same uniqueness, which proves \textup{(iv)}.
\end{proof}

\subsection{Complex parameter disks and the growth estimate}
\label{model:growth-subsection}

The real domains in \eqref{model:real-domains} have a cusp at
\(T=0\).  We fill graphs over a fixed parameter half-disk to
produce a parameter disk of radius comparable to \(D^{1/r}\)
at a phase point of depth \(D\).  The shape constraint on this
disk will bound the second parameter coefficient.

\begin{lemma}[Holomorphic filling of a moving graph]
\label{model:graph-filling}
Let \(\Gamma\) be a piecewise smooth Jordan curve, and let \(U\)
be a connected open neighborhood of a compact path in a complex
variable \(w\).  Suppose \(A(T,w)\) is holomorphic on a
neighborhood of \(\Gamma\times U\).
If, for \(w\) in some nonempty open subset of \(U\), it extends
holomorphically to the interior of \(\Gamma\), then it has such
an extension for every \(w\in U\), jointly holomorphic in \(T,w\).
The extension agrees with the original collar germs.
\end{lemma}

\begin{proof}
For every integer \(n\ge0\), the moment
\[
 \int_\Gamma T^n A(T,w)\,\dd T
\]
is holomorphic in \(w\).  It vanishes in the initial open set
and hence throughout \(U\).  For an exterior point \(z\) of
large modulus, expand \((T-z)^{-1}\) in its geometric series.
The vanishing moments show that the exterior Cauchy transform
is zero there, and analytic continuation makes it zero on the
whole exterior.  The interior Cauchy integral
\[
 \frac{1}{2\pi\ii}\int_\Gamma
       \frac{A(T,w)}{T-z}\,\dd T
\]
is jointly holomorphic in \(z,w\).  Deforming the contour a
short distance inside and outside any collar point and using
the zero exterior transform shows that it agrees there with
\(A\).  It is the required extension.
\end{proof}

\begin{proposition}[Growth of the limiting first field]
\label{prop:model-growth}
For every compact horizontal interval \(I\subset\R\) and every
\(\eps>0\), there is a constant \(C_{I,\eps}\) such that
\begin{equation}
\label{model:power-growth}
 |F(x-\ii D)|\le C_{I,\eps}D^{-\lambda-\eps}
 \qquad(x\in I,\ 0<D<1).
\end{equation}
Equivalently,
\begin{equation}
\label{model:log-growth}
 \sup_{x\in I}\log|F(x-\ii e^{-v})|
       \le\lambda v+o(v)\qquad(v\to\infty),
\end{equation}
where the error is uniform on \(I\).
\end{proposition}

\begin{proof}
Fix \(\eps>0\), use \Cref{prop:pure-model} with
\(r=1/(\lambda+\eps)\in(0,1)\), and take a small fixed upper
half-disk \(|T|<t_0\), \(\Im T>0\).
For a target \(Y\) with \(\Re Y\in I\), \(-\Im Y=D\),
and for \(|T_c|\le c_*D^{1/r}\) in this half-disk, use the
graph with phase argument
\begin{equation}
\label{model:parameter-graph}
 Y+G(T)-G(T_c),\qquad
 G(T)=-\ii C_*(\eta_0-\ii T)^r,\qquad
 0<\eta_0\le c_*D^{1/r}.
\end{equation}
The fractional power is the branch on the right half-plane;
\(\Re(\eta_0-\ii T)>0\) in the upper half-disk.
The auxiliary \(\eta_0\) is distinct from the fixed
growth exponent \(\eps\).

On the real diameter,
\[
 -\Im G(T)=C_*\Re(\eta_0-\ii T)^r
       \ge C_*\cos(\pi r/2)|T|^r.
\]
Choose \(C_*\) large compared with the constant in
\eqref{model:real-domains}, and then \(c_*\) so small that
\(|G(T_c)|<D/2\).  Every diameter point then lies strictly
in its real-parameter domain.  On the semicircle
\(|T|=t_0\), \(\Im T\ge0\),
\[
 -\Im G(T)\ge c_rC_*t_0^r,\qquad
 |\Re G(T)|\le C_r(-\Im G(T)).
\]
Fix the required cone aperture first and then increase
\(C_*\) further to put this semicircle in the deep joint
domain.  The bounded interval \(I\) and the correction
\(G(T_c)\) cost only fixed additional margins.

The graph has coherent joint germs on a boundary collar.
On the diameter these are the germs of
\Cref{prop:pure-model}\textup{(iv)}; on the semicircle they
are deep joint data.  They agree at the corners because
both are attached to the same deep branch.  Overlapping
real-parameter neighborhoods agree first for real parameters
and then by holomorphic identity.

Lower the center \(Y\) until the whole graph lies in deep
joint data, then raise it vertically to the target.
Boundary membership and strict margins persist.  On each
compact center segment a finite covering gives a holomorphic
function on a neighborhood of the boundary contour times
a connected complex neighborhood of that segment.
Coherence follows from vertical-path uniqueness.  Initially
the graph interior is filled by deep data, so
\Cref{model:graph-filling} fills it throughout the motion.
Its boundary values have negative imaginary part; the
harmonic maximum principle therefore gives \(\Im S<0\)
on the filled interior.

For the lower half-disk use
\(G(T)=-\ii C_*(\eta_0+\ii T)^r\).
The two constructions agree on the joint real-parameter
germs.  The germs obtained for different \(T_c\) also agree
on overlaps: at a fixed parameter they follow the same
vertical path from the unique deep branch.  Thus the
model is holomorphic on the whole parameter disk
\begin{equation}
\label{model:small-parameter-disk}
 |T|<c_*D^{1/r}
\end{equation}
at the target, still with \(\Im S<0\).
Shrinking \(c_*\) once makes this construction uniform
for \(Y\) in a disk of radius a small fixed multiple of \(D\).

Write
\[
 S(T,Y)=Y+T^2A(Y)+O(T^4).
\]
Evenness follows from the deep branch and continuation.
The positive harmonic function \(-\Im S(T,Y)\) equals
\(D\) at \(T=0\).  Its Fourier coefficient estimate on
circles in \eqref{model:small-parameter-disk}, or the
Caratheodory coefficient estimate, gives
\[
 |A(Y)|\le C D^{1-2/r}.
\]
This also holds on the depth-sized \(Y\)-neighborhood,
so Cauchy's estimate yields
\begin{equation}
\label{model:A-derivative-bound}
 |A'(Y)|\le C D^{-2/r}.
\end{equation}
The energy equation determines the derivative:
\[
 g(z)=1-\frac{z^2}{10}+O(z^4),\quad
 J=1+T^2A'(Y)+O(T^4),\quad
 z=\frac{TF(Y)}2+O(T^3).
\]
Its coefficient of \(T^2\) is
\(3A'(Y)-F(Y)^2/40=0\).  Hence
\begin{equation}
\label{model:A-coefficient}
 A'(Y)=F(Y)^2/120.
\end{equation}
Combining these bounds gives
\(|F(Y)|\le C D^{-1/r}=C D^{-\lambda-\eps}\).
Arbitrariness of \(\eps>0\) gives
\eqref{model:log-growth}.
\end{proof}

\subsection{Uniform high-value witnesses}
\label{model:witness-subsection}

Set
\begin{equation}
\label{model:L-definition}
 L(v)=\log m_*(e^{-v}).
\end{equation}
By \eqref{model:star-moment},
\(L(v)=\lambda v+o(|v|)\) at both ends.

\begin{lemma}[Common normal epochs]
\label{model:normal-epochs}
Fix
\[
 \frac12<c_1<1-\beta,\qquad
 P_1>\max(\lambda,1+a).
\]
There is an unbounded set \(E\subset(0,\infty)\), with
\(\dist(v,E)=o(v)\) as \(v\to\infty\), such that every
\(u\in E\) satisfies, for all real \(v\),
\begin{equation}
\label{model:epoch-envelope}
 L(v)-L(u)\le
 \begin{cases}
  c_1(v-u),&v\le u,\\
  P_1(v-u),&v\ge u.
 \end{cases}
\end{equation}
\end{lemma}

\begin{proof}
The coarse part of \eqref{model:new-global-envelope} gives
\(L(v)\le C+(1-\beta)v\) for \(v\le0\).
Thus \(L(v)-c_1v\to-\infty\) as \(v\to-\infty\).
For a large \(w\), choose a running record \(u_0\le w\)
of \(L(v)-c_1v\) on \((-\infty,w]\).
Its positive asymptotic slope \(\lambda-c_1\) implies
\(w-u_0=o(w)\): for every fixed \(\eta>0\), all values
at \(v\le(1-\eta)w\) are eventually below the value at \(w\).

Next maximize \(L(v)-P_1v\) over \(v\ge u_0\), and call
a maximizer \(u\).  It exists because \(P_1>\lambda\).
The negative asymptotic slope shows \(u-u_0=o(u_0)\).
Its maximality gives the future bound in
\eqref{model:epoch-envelope}.
For \(v\le u_0\), the past record and
\(L(u)-L(u_0)\ge P_1(u-u_0)\ge c_1(u-u_0)\)
give the past bound at \(u\).  For \(u_0\le v\le u\),
maximality gives
\[
 L(v)-L(u)\le P_1(v-u)\le c_1(v-u).
\]
The epochs obtained this way are within \(o(w)\) of
every sufficiently large \(w\).
\end{proof}

\begin{proposition}[Witnesses near every horizontal center]
\label{prop:model-witnesses}
There is a fixed sufficiently large aperture \(L_0\) with
the following property.  For every compact interval
\(I\subset\R\), there is an error function
\(\eta_I(u)\to0\), for \(u\in E\), such that every
sufficiently large \(u\in E\) and every \(x\in I\) admit
a point \(w=p-\ii\delta\in\mathbb H_-\) satisfying
\begin{equation}
\label{model:witness-estimate}
 |p-x|\le L_0\delta,\qquad
 |-\log\delta-u|\le\eta_I(u)u,\qquad
 \log|F(w)|\ge\lambda u-\eta_I(u)u.
\end{equation}
The epochs are common to all centers; a center-dependent
subsequence is not needed.
\end{proposition}

\begin{proof}
Fix
\[
 \frac12<\alpha'<\min(\alpha,c_1,1)
\]
and choose \(L_0\) large enough for the quiet-cone form
of \Cref{thm:no-trivial-profile} with this exponent.
If the conclusion fails, some fixed small \(\eta>0\),
epochs \(u_j\to\infty\) in \(E\), and centers \(x_j\)
in a fixed compact interval \(I\) satisfy
\begin{equation}
\label{model:failed-witness}
 \log|F(p-\ii e^{-v})|<\lambda u_j-\eta u_j
\end{equation}
for every
\[
 |p-x_j|\le L_0e^{-v},\qquad |v-u_j|\le\eta u_j.
\]
Pass to a subsequence with \(u_j\ge j^2\).

Return to the actual data of
\Cref{model:selected-fields}, taking their sequence index
sufficiently large after each fixed \(u=u_j\).  Set
\begin{equation}
\label{model:witness-rescaling}
 d''=d'e^{-u},\qquad K''=K+u,\qquad B''=\mu(K+u),
\end{equation}
with physical center \(\xi_n+d'x_j\).
We may require \(K'',B''\to\infty\): for fixed \(u\),
\(B''/B'\) tends to the positive number
\(e^{-u}m_*(e^{-u})\).

First retain a uniform global moment envelope at these
scales.  Write
\[
 M_n(v)=e^v\frac{\mu(K+v)}{B'},\qquad
 M_n(v)\longrightarrow e^{L(v)}
\]
locally uniformly.  The old envelope says
\begin{equation}
\label{model:witness-old-tail}
 M_n(v)\le
 \begin{cases}
  C e^{(1-\beta)v},&v\le0,\\
  C e^{(1+a)v},&v\ge0.
 \end{cases}
\end{equation}
Choose fixed \(1/2<c_2<c_1\) and
\(P_2>\max(P_1,1+a)\).
On the finite range \(0\le v\le C_0u\), take the
actual index so large that the ratios
\(M_n(v)/M_n(u)\) differ from their limits by at most
a factor two.  The epoch bounds imply
\begin{align}
\label{model:witness-local-envelope}
 \frac{M_n(v)}{M_n(u)}
 &\le2e^{c_2(v-u)}&& (0\le v\le u),\\
\label{model:witness-local-fine}
 \frac{M_n(v)}{M_n(u)}
 &\le2e^{P_2(v-u)}&& (u\le v\le C_0u).
\end{align}
For all \(v\le0\), \eqref{model:witness-old-tail}
supplies the first bound up to one fixed constant,
since \(1-\beta>c_2\) and
\(L(u)=\lambda u+o(u)\), \(\lambda\ge1>c_2\).
For all \(v\ge C_0u\), it supplies the second bound
if \(C_0\) is a sufficiently large fixed constant:
\[
 (1+a-P_2)v+P_2u-L(u)\le0.
\]
In terms of \(\mu\), this is a uniform new envelope
with coarse exponent \(1-c_2<1/2\) and fine exponent
\(P_2-1>0\).  Only a finite range chosen before the
actual index used moment convergence.

We next prove quietness on the whole needed cone,
including the unbounded coarser tail.  At old log
depth \(v\), the denominator for the new weighted
quiet-cone assertion is
\begin{equation}
\label{model:witness-denominator}
 M_n(u)e^{\alpha'(v-u)}.
\end{equation}
Indeed the new depth ratio is \(e^{u-v}\), and
\[
 \frac{d''}{B''}M_{\rm shape}
 =
 \frac{\max(|F_n|,\sqrt{|Q_n|})}{M_n(u)},
\]
where \(M_{\rm shape}=\max(|f|,\sqrt{|q|})\).

Increase the actual index further so that
\begin{equation}
\label{model:witness-diagonal}
 |F_n-F|\le1,\qquad \sqrt{|Q_n|}\le1
\end{equation}
throughout
\[
 0\le v\le u_j+j,\qquad |p-x_j|\le L_0e^{-v}.
\]
For each \(j\), a small enlargement of this set is a
compact subset of \(\mathbb H_-\), so compact convergence
allows the choice.  Throughout the set,
\eqref{model:witness-denominator} is at least
\[
 \exp\bigl((\lambda-\alpha')u_j+o(u_j)\bigr).
\]
The errors in \eqref{model:witness-diagonal} are
therefore negligible uniformly.

For \(u_j-\eta u_j\le v\le u_j+j\), the upper endpoint
is in the failed-witness window for large \(j\).
Equation \eqref{model:failed-witness} bounds the
logarithmic ratio by
\[
 -(1-\alpha')\eta u_j+o(u_j).
\]
For \(0\le v<u_j-\eta u_j\), use
\eqref{model:log-growth} on the fixed horizontal
enlargement of \(I\) containing this part of the cone.
Its errors over \(0\le v\le u_j\) are \(o(u_j)\):
separate a fixed compact \(v\)-range, then use the
uniform asymptotic estimate on the remainder.
The logarithmic ratio is at most
\[
 (\lambda-\alpha')(v-u_j)+o(u_j)
 \le-(\lambda-\alpha')\eta u_j+o(u_j).
\]

For every \(v<0\) back to fixed physical depth,
no convergence transfer is necessary.
If \(I\subset[-R,R]\), the cone at old depth
\(e^{-v}\ge1\) lies inside the old cone of aperture
\(L_0+R\), because
\[
 |p|\le R+L_0e^{-v}\le(L_0+R)e^{-v}.
\]
The inherited bound \eqref{model:selected-cones} gives
\[
 \frac{\max(|F_n|,\sqrt{|Q_n|})}
      {M_n(u_j)e^{\alpha'(v-u_j)}}
 \le
 C_{L_0+R}
 \exp\bigl(-(\lambda-\alpha')u_j+o(u_j)\bigr)
 e^{(\alpha-\alpha')v}.
\]
It tends uniformly to zero since \(\alpha>\alpha'\).
Both the aperture and its constant remain fixed
despite movement of the centers.

The three ranges give quietness down to \(e^{-j}d''\),
and hence down to every fixed positive multiple of
\(d''\) eventually.  The normalization is exactly
\(B''=\mu(K'')\), with the uniform moment envelope
proved above.  These are the hypotheses of the
quiet-cone form of \Cref{thm:no-trivial-profile}.
For clarity, at a fixed depth ratio and bounded
scaled parameter, the parameter-to-depth ratio
tends to zero because \(B''\to\infty\).
The local cone theorem uses a fixed ratio bound and
a fixed aperture.  Its quiet passage forces the
global real width deficit divided by \(d''\) to
vanish, and the one-sided coefficient and moment
argument in that theorem is contradictory.

Thus, for each fixed small \(\eta>0\), all sufficiently
large common epochs have the required witnesses
uniformly over \(x\in I\).  A decreasing diagonal
choice of \(\eta\) yields \(\eta_I(u)\to0\) in
\eqref{model:witness-estimate}.
\end{proof}

\subsection{The high-modulus mesh contradiction}

We isolate the function-theoretic consequence of the growth bound and the
common-epoch witnesses.  In the following argument a depth is the positive
number \(\delta=-\Im z\), and its logarithmic depth is
\(v=-\log\delta\).

\begin{theorem}[High-modulus mesh contradiction]
\label{thm:mesh-contradiction}
Let \(F\not\equiv0\) be holomorphic in
\(\mathbb H_-:=\{z\in\C:\Im z<0\}\), and suppose that \(F''/F\)
extends holomorphically to \(\mathbb H_-\).  There cannot exist a number
\(\lambda\ge1\) and an unbounded set of epochs
\(E\subset(0,\infty)\) with all of the following properties.
\begin{enumerate}[label=\textup{(\roman*)}]
\item For every compact interval \(I\subset\R\),
\begin{equation}
 \sup_{x\in I}\log|F(x-\ii e^{-v})|
 \le \lambda v+o(v),\qquad v\longrightarrow\infty.
 \label{model:mesh-upper}
\end{equation}
The error is uniform in \(x\in I\).
\item The epochs are relatively dense in the sense that
\begin{equation}
 \dist(v,E)=o(v),\qquad v\longrightarrow\infty.
 \label{model:mesh-epochs}
\end{equation}
\item For every compact interval \(I\), there are a fixed aperture
\(A_I<\infty\) and a function \(r_I(u)\to0\) such that, for every
sufficiently large \(u\in E\) and every \(x\in I\), there is a point
\(w=p-\ii\delta\) satisfying
\begin{equation}
 |p-x|\le A_I\delta,\qquad
 |{-\log\delta}-u|\le r_I(u)u,\qquad
 \log|F(w)|\ge\lambda u-r_I(u)u.
 \label{model:mesh-witness}
\end{equation}
\end{enumerate}
In particular, the epochs in \textup{(iii)} are common to all the centers
in \(I\); no choice of an epoch depending on the center is allowed.
\end{theorem}

We first prove the two local estimates used to construct the mesh.
All references to Brownian motion below can equivalently be formulated
using harmonic measure and the minimum principle for superharmonic
functions.

\begin{lemma}[Small shadows of a high-deficit set]
\label{model:mesh-shadow}
Fix constants \(1\le a\le b\) and \(K>0\).  Let \(w=p-\ii d\), and let
\(\mathcal T\) be a finite family of squares \(T\) of sides \(\rho_T\)
and center depths \(d_T\), with
\[
 a\rho_T\le d_T\le b\rho_T,\qquad
 d_T\le d/2,\qquad |\Re c_T-p|\le Kd.
\]
Choose a rectangle
\[
 R=\{z:|\Re z-p|<Md,\quad \tau<-\Im z<2d\},
\]
where \(M\) is sufficiently large in terms of \(a,b,K\), and
\(0<\tau<\frac14\min_T d_T\).  Suppose that \(U\ge0\) is
superharmonic in a neighborhood of \(\overline R\), that \(U(w)<\infty\),
and that \(U\ge H>0\) on the concentric disks \(D_T\) of radius
\(\rho_T/8\).  If \(I_T\) is the horizontal projection of \(T\), then
\begin{equation}
 \left|\bigcup_{T\in\mathcal T}I_T\right|
 \le C d\,\frac{U(w)}{H}.
 \label{model:mesh-shadow-bound}
\end{equation}
Here \(C\) depends only on the fixed geometric constants.
\end{lemma}

\begin{proof}
Let \(B\) be Brownian motion started at \(w\), stopped when it exits
\(R\).  Superharmonic comparison, applied first up to the first hit of
the union of the disks and then up to the exit, gives
\begin{equation}
 \Prob_w\{B\text{ hits some }D_T\text{ before exiting }R\}
 \le U(w)/H.
 \label{model:mesh-hitting-upper}
\end{equation}
This comparison allows logarithmic poles of \(U\): one can first stop
away from those poles and then use monotone exhaustion.  Alternatively,
the harmonic probability of hitting the disk union, multiplied by
\(H\), is a subfunction of \(U\) by the minimum principle on the
complement of the disks.

For each \(T\), consider the event that the motion hits \(D_T\) and
subsequently exits \(R\) through its top side at a horizontal coordinate
in \(I_T\).  We claim that its probability is at least
\(c\rho_T/d\).  First stop at the horizontal line of depth \(d_T\), or
at the other sides of the part of \(R\) below that line.  The harmonic
measure density on the segment
\[
 |x-\Re c_T|<\rho_T/32,
 \qquad -\Im z=d_T,
\]
is bounded below by \(c/d\).  To see the uniformity, divide all lengths
by \(d\).  The starting point has depth one, the bottom has depth two,
and the new top has depth in \([0,1/2]\); the segment lies a fixed
positive distance from the lateral sides.  The interior-to-side
Poisson density of a rectangle is positive and continuous in these
parameters, including top depth zero.  Its minimum on this compact
family is therefore positive.  Thus the probability of this first hit
inside \(D_T\) is at least \(c\rho_T/d\).

From any point of that segment, the probability of exiting the original
rectangle through its top within \(I_T\) has a fixed positive lower
bound.  Indeed, use a smaller rectangle with horizontal center
\(\Re c_T\), horizontal half-width a fixed large multiple of \(d_T\),
top depth \(\tau\), and bottom depth \(2d_T\).  It lies in \(R\) after
increasing \(M\).  After division by \(d_T\), its starting points,
top depth in \([0,1/4]\), and target intervals of length
\(\rho_T/d_T\in[1/b,1/a]\) again range over a compact family with
nondegenerate margins.  Positivity and continuity of harmonic measure
give the claimed lower bound.  Exiting this smaller rectangle on the
specified top interval also exits \(R\) there.  The strong Markov
property proves the claim.

Choose a subfamily of the intervals \(I_T\) with disjoint interiors
whose total length is at least one third of the length of their union.
For completeness, order the finite family by decreasing length, retain
an interval if its interior is disjoint from those already retained,
and discard it otherwise.  Each discarded interval is contained in the
threefold concentric enlargement of an intersecting retained interval
of at least its length.  This proves the asserted covering estimate.
The events just constructed for the retained intervals are disjoint,
because their final exit coordinates are in disjoint intervals;
individual endpoints have harmonic measure zero.  Each event is
contained in the event in \eqref{model:mesh-hitting-upper}.  Summing their
lower bounds gives \eqref{model:mesh-shadow-bound}.  In particular, no
union bound over the number of tiles is being used.
\end{proof}

\begin{lemma}[Routing near a good square]
\label{model:mesh-routing}
Let \(Q\) be a square of side \(\rho\) and center \(c\).  Suppose that
\(F\) is holomorphic on \(B(c,100\rho)\), and that
\[
 U=M-\log|F|\ge0
 \quad\text{there},\qquad
 \inf_Q U\le T,
\]
where \(T>0\).  There are disks with sum of radii at most
\(C h\rho\), for any fixed sufficiently small \(h>0\), outside which
\begin{equation}
 U(z)\le C_hT
 \quad\text{on the concentric square }Q^*\text{ of side }9\rho.
 \label{model:mesh-local-high}
\end{equation}
The constants are absolute except for the indicated dependence on
\(h\).  A finite simple chain of such squares, in which touching
neighbors have side ratios in \([1/2,2]\), can consequently be routed
through their enlargements by a polygonal path satisfying the same
bound with length at most \(C\sum_Q\rho_Q\).  The endpoints can be
chosen in prescribed squares of side comparable to \(\rho_Q\)
near the first and last tiles and contained in their respective
enlargements; the routing constants may depend on these fixed
comparability constants.
\end{lemma}

\begin{proof}
Choose \(q\in Q\) with \(U(q)\le2T\).  A disk about \(q\) of radius
\(40\rho\) lies in the given domain and contains \(B(c,30\rho)\).
The super-mean inequality, integrated over radii, and nonnegativity give
\begin{equation}
 \int_{B(c,30\rho)}U\,\dd A\le C\rho^2T.
 \label{model:mesh-local-L1}
\end{equation}
The distributional identity
\(-\Delta U=2\pi\sum_a m_a\delta_a\) records the zeros of \(F\)
with their multiplicities.  A nonnegative smooth cutoff equal to one
on \(B(c,20\rho)\), supported in \(B(c,30\rho)\), and with second
derivatives bounded by \(C\rho^{-2}\) yields
\[
 N:=\sum_{a\in B(c,20\rho)}m_a
 \le C\rho^{-2}\int_{B(c,30\rho)}U\,\dd A\le CT.
\]
List these zeros with multiplicity as \(a_1,\ldots,a_N\).  The function
\[
 H(z)=U(z)+\sum_{\nu=1}^N
               \log\frac{|z-a_\nu|}{\rho}
\]
extends harmonically across the zeros on \(B(c,20\rho)\).  Each
logarithm has bounded absolute integral in scaled coordinates, so
\eqref{model:mesh-local-L1} gives
\(\int_{B(c,20\rho)}|H|\,\dd A\le C\rho^2T\).
The mean-value inequality for the absolute value of a harmonic function
then gives \(|H|\le CT\) on \(B(c,10\rho)\).

Here is the precise elementary covering estimate needed for the
logarithmic terms.  For points \(b_1,\ldots,b_N\in\C\), outside disks
of sum of radii at most \(3h\),
\begin{equation}
 \prod_{\nu=1}^N|z-b_\nu|\ge(h/e)^N.
 \label{model:mesh-cartan}
\end{equation}
There is nothing to prove when \(N=0\).  Otherwise, consider all disks
\(B(z,hi/N)\) whose interiors contain at least \(i\) of the listed
points, for \(1\le i\le N\).  Starting with the largest of the finitely
many possible radii, select a maximal disjoint family at each radius,
discarding disks that intersect ones already selected.  Only finitely
many disks can be selected: their centers lie in a bounded set and their
radii are at least \(h/N\).  Each unselected disk intersects a selected
one of at least its radius, and is contained in that disk's threefold
enlargement.  If the selected radii are \(hi_k/N\), disjoint interiors
show that \(\sum_k i_k\le N\), by counting the points they contain.
The enlarged disks therefore have total radius at most \(3h\).
Outside them, the \(i\)-th smallest distance to the points is at least
\(hi/N\), for every \(i\).  Their product is at least
\(h^N N!/N^N\ge(h/e)^N\), proving \eqref{model:mesh-cartan}.

Apply this estimate to \(b_\nu=(a_\nu-c)/\rho\).  Outside the
resulting disks,
\[
 U\le CT+N\log(e/h)\le C_hT
 \quad\text{on }B(c,10\rho).
\]
This implies \eqref{model:mesh-local-high}.

It remains to justify actual routing, including at corner contacts.
If \(Q,Q'\) touch and have comparable sizes, choose a square centered
at a contact point of side \(\min(\rho_Q,\rho_{Q'})\).  This square
lies in \(Q^*\cap(Q')^*\).  The sums of the lengths of the horizontal
and vertical projections of all exceptional disks for both tiles are
at most \(C h(\rho_Q+\rho_{Q'})\).  Fix \(h\) so small that each of
these sums is less than half the side of the contact square.  One can
then choose an \(x\)-coordinate and a \(y\)-coordinate in that square,
each outside the corresponding projections of both disk families.
Their intersection is a junction point safe for both tiles.  Choose
the two endpoint coordinates in the same fashion in the prescribed
endpoint squares.  Within a tile, join its two chosen points by one
horizontal and one vertical segment.  Both coordinates at each endpoint
avoid all exceptional projections for that tile, so these segments
avoid all of its exceptional disks.  They lie in \(Q^*\), and their
combined length is at most \(C\rho_Q\).  Concatenation proves the
routing assertion.  The same construction works when a chain consists
of a single tile.
\end{proof}

\begin{proof}[Proof of \Cref{thm:mesh-contradiction}]
We work in fixed horizontal intervals contained in \([-5,5]\).  All
errors below are uniform on these intervals.  Fix the aperture from
\eqref{model:mesh-witness} for this compact interval.

\paragraph{Separated bands with neighboring logarithmic depths.}
Choose a nonincreasing function \(\theta(v)\to0\), with
\(\theta(v)\ge v^{-1/4}\), whose square bounds, eventually, the
relative errors in \eqref{model:mesh-epochs} and
\eqref{model:mesh-witness}, uniformly for all arguments at least \(v\).
Such a majorant is obtained by taking the decreasing tail supremum of
these errors, its square root, and the maximum with \(v^{-1/4}\).
Enlarging it slightly accommodates approximate nearest epochs if
\(E\) is not closed.

Starting at a sufficiently large \(u_1\in E\), set
\(s_j=\theta(u_j)u_j\), and choose
\begin{equation}
 u_{j+1}\in E,\qquad
 u_{j+1}=u_j+8s_j+o(s_j).
 \label{model:mesh-band-spacing}
\end{equation}
Indeed the error in finding an epoch near \(u_j+8s_j\) is at most
\(C\theta(u_j)^2u_j=o(s_j)\).
The monotonicity of \(\theta\) implies
\(s_{j+1}\le(1+o(1))s_j\).  Consequently the bands
\([u_j-s_j,u_j+s_j]\) are disjoint, their gaps tend to infinity,
\(s_j\to\infty\), and
\begin{equation}
 u_{j+k}=u_j+o(u_j)\qquad (k=-1,0,1,2).
 \label{model:mesh-neighboring-scales}
\end{equation}
Every witness at epoch \(u_j\) has logarithmic depth
\(u_j+o(s_j)\).  These statements survive any fixed additive changes
to the ends of the logarithmic bands.

Choose dyadic depths \(a_j,b_j\) with
\[
 a_j=e^{-u_j-s_j+O(1)},\qquad
 b_j=e^{-u_j+s_j+O(1)},\qquad a_j<b_j.
\]
The horizontal rectangle is
\(R_j=[-3,3]\times[a_j,b_j]\) in the coordinates \((x,\delta)\).
The two vertical rectangles have horizontal intervals
\([-2,-3/2]\) and \([3/2,2]\), and depth interval
\([a_{j+2}/8,8b_j]\); dyadic rounding can be incorporated into these
fixed factors.  The factors eight leave room at the ends for all
subsequent local enlargements.  A witness at epoch \(u_{j-1}\) is
deeper than the entire vertical rectangle by a factor tending to
infinity.  In fact, the logarithmic gap from that witness to its bottom
is at least \((7-o(1))s_{j-1}\).

Take a slightly larger open region containing these three rectangles,
the witnesses at epochs \(u_j,u_{j-1}\) needed for centers in
\([-3,3]\), and the fixed-factor boxes that will be used about them.
Its horizontal coordinates lie in
\([-5,5]\); all its logarithmic depths are
\(u_j+o(u_j)\), by \eqref{model:mesh-neighboring-scales}.  By
\eqref{model:mesh-upper} we can choose
\begin{equation}
 M_j=\lambda u_j+o(u_j),\qquad M_j\ge\lambda u_j,
 \qquad U_j=M_j-\log|F|\ge0
 \label{model:mesh-deficit}
\end{equation}
throughout this region.  The function \(U_j\) is superharmonic and may
have logarithmic poles at the zeros of \(F\).  Its value at every
witness just specified is \(o(u_j)\), uniformly in its center.

\paragraph{Tiles and the exclusion of bad barriers.}
Fix a sufficiently large dyadic integer \(N_0\).  In the layer
\(2^{-l-1}\le\delta\le2^{-l}\), tile by squares of side
\(\rho=2^{-l-1}/N_0\), using the common dyadic horizontal grid.
The sides of a tile are thus a small fixed fraction of its depth.
Take \(N_0\) large enough that the disks of radius \(100\rho\)
required in \Cref{model:mesh-routing} lie within the fixed-factor
enlargements already allowed.  The rectangles have dyadic boundaries,
so they are finite unions of these tiles.  Tiles that touch, even at a
corner, have side ratios between one half and two.

Choose \(\eps_j\to0\) so slowly that the uniform witness deficits in
\eqref{model:mesh-deficit} are \(o(\eps_j u_j)\), and so that
\(M_j-\lambda u_j=o(\eps_j u_j)\).  Color a tile good when
\begin{equation}
 \sup_T\log|F|\ge\lambda u_j-\eps_j u_j,
 \label{model:mesh-good-tile}
\end{equation}
and bad otherwise.  On every bad tile, \(U_j\ge\eps_j u_j\).
For any witness \(w\) at one of the two indicated epochs,
\Cref{model:mesh-shadow} therefore implies that the union of the
projections of bad tiles with center depth at most \(\delta_w/2\)
and horizontal center distance at most \(K\delta_w\) from \(w\)
has length
\begin{equation}
 o(\delta_w).
 \label{model:mesh-small-bad-shadow}
\end{equation}
Here \(K\) is any fixed constant needed below.  For this application
put the fictitious top of the Brownian rectangle a fixed factor
shallower than the shallowest tile under consideration.  The rectangle
then lies in the enlarged region of \eqref{model:mesh-deficit}.

There is no bad chain in \(R_j\) from its shallow edge to its deep
edge, even when corner contacts are allowed.  Suppose such a chain
exists and let \(x_0\) be a point where it meets the shallow edge.
Choose the epoch-\(u_j\) witness \(p-\ii\delta_w\) for the center
\(x_0\).  The shallow edge is much shallower than \(\delta_w\), and
the deep edge is much deeper.  Follow the chain from its shallow end
until it first reaches depth \(\delta_w/4\) or first leaves the
horizontal interval \([x_0-\delta_w,x_0+\delta_w]\).  Taking \(N_0\)
large guarantees that all tiles up to and including this first event
have center depths at most \(\delta_w/2\) and center distances from
\(p\) at most \((A_I+2)\delta_w\).  In the first case the last tile
has side at least \(c\delta_w/N_0\), so its shadow alone has length
at least a fixed positive multiple of \(\delta_w\).  In the second
case the projections of this connected chain prefix cover an interval
of length at least \(c\delta_w\).  Both conclusions contradict
\eqref{model:mesh-small-bad-shadow}.

Neither vertical rectangle has a bad chain joining its two lateral
sides.  Choose an epoch-\(u_{j-1}\) witness for the midpoint of its
fixed horizontal interval.  This witness has depth \(\delta_w\),
and every tile of the rectangle has center depth less than
\(\delta_w/2\) for large \(j\).  Its horizontal coordinate \(p\)
lies in the interior of the strip, since its distance from the
midpoint is at most \(A_I\delta_w\to0\).  A lateral bad crossing
projects onto the whole strip.  The tiles of that crossing with
centers in \([p-2\delta_w,p+2\delta_w]\) therefore have projections
covering \([p-\delta_w,p+\delta_w]\), once \(N_0\) is large.  This
again contradicts \eqref{model:mesh-small-bad-shadow}.

\paragraph{Planar duality and routing the good chains.}
We include the finite grid argument to specify the corner convention.
Subdivide every tile of a rectangle to its finest dyadic square grid,
retaining the parent tile's color.  On a finite square grid, either
there is a good left-to-right path using shared edges, or there is a
bad top-to-bottom path using shared edges or corners.  To prove this,
mark all good cells reachable by shared edges from the left side and
adjoin an exterior collar along the left.  If this set reaches the
right, the first alternative holds.  Otherwise follow its right-hand
boundary from the top of the rectangle to the bottom, with the marked
set on the left.  At vertices with several possible continuations,
keep the marked set on the left, equivalently resolve the vertex by a
small rounded arc.  Closed boundary components can be ignored.  The
cells immediately on the right of the resulting boundary are bad:
an edge-adjacent good cell would have been marked.  Consecutive such
cells share an edge or a corner, and they join the top to the bottom.
If no cell was marked, the first column itself is such a bad chain.
This proves the alternative.  Interchanging the two coordinate
directions gives its vertical version.  Mapping a chain back to its
parent tiles preserves corner adjacency; removing loops and repeated
tiles gives a simple parent chain.

The bad barriers have already been excluded.  Thus \(R_j\) has a good
horizontal chain, and each vertical rectangle has a good vertical
chain.  On a good tile the infimum of \(U_j\) is at most
\(M_j-\lambda u_j+\eps_j u_j=o(u_j)\), uniformly.  Apply
\Cref{model:mesh-routing} to route each chain.  We obtain a horizontal
path \(\gamma_j\), extending from near \(-3\) to near \(3\), and
vertical paths \(\sigma_j^-\), \(\sigma_j^+\) in slight enlargements
of the two side strips, with
\begin{equation}
 \log|F|\ge\lambda u_j-o(u_j)
 \quad\text{on }\gamma_j\cup\sigma_j^-\cup\sigma_j^+.
 \label{model:mesh-path-height}
\end{equation}
The horizontal path stays in a fixed-factor enlargement of its band.
The vertical paths start deeper than all of \(\gamma_j\) and end
shallower than all of \(\gamma_{j+2}\).  They remain to the left or
right, respectively, of the entire interval \([-1,1]\).

The length estimate in the routing lemma is subexponential in \(u_j\).
Indeed, in any of these rectangles the smallest tile side is
\(e^{-u_j+o(u_j)}\), the largest side is
\(e^{-u_j+o(u_j)}\), its height is \(e^{-u_j+o(u_j)}\), and its
width is bounded.  Even its entire finest grid contains at most
\(e^{u_j+o(u_j)}\) cells.  Hence the number of parent tiles times
their largest side, and therefore each routed path's length, is at
most
\begin{equation}
 \exp(o(u_j)).
 \label{model:mesh-path-length}
\end{equation}
All local constants, including those in the Cartan covering, have
been fixed before \(j\) tends to infinity.

Each vertical path intersects each of
\(\gamma_j,\gamma_{j+1},\gamma_{j+2}\).  One elementary way to
verify this is to erase loops from a horizontal polygonal path and
trim it at two vertical lines outside the side strips.  It is then a
simple arc joining the opposite sides of a rectangle, confined to its
horizontal band.  Closing this arc along three sides of the rectangle
and applying the Jordan curve theorem separates points above the
band from points below it.  A vertical path in either side strip has
endpoints on these different sides and must meet the arc.  The
fixed-factor margins ensure that local routing has not changed this
separation.

\paragraph{A primitive with no logarithmic monodromy.}
Set \(q=F''/F\), holomorphic by hypothesis.  Every zero \(z_0\) of
\(F\) is simple.  Otherwise \(F(z_0)=F'(z_0)=0\), and uniqueness
for the holomorphic equation \(F''=qF\) would imply \(F\equiv0\).
At a simple zero this equation also gives \(F''(z_0)=0\), and hence
\[
 F(z)=a(z-z_0)+O((z-z_0)^3),\qquad
 F(z)^{-2}=a^{-2}(z-z_0)^{-2}+O(1),\quad a\ne0.
\]
In particular all residues of \(F^{-2}\) vanish.  Integration around
any closed curve avoiding its poles is zero: the curve is compact in
the simply connected half-plane and encloses only finitely many
poles, so this is the residue theorem.  Thus there is a single-valued
meromorphic function \(C_0\) in \(\mathbb H_-\) with
\begin{equation}
 C_0'=F^{-2}.
 \label{model:mesh-primitive}
\end{equation}
Its singularity at a zero \(z_0\) is
\(-a^{-2}(z-z_0)^{-1}\).

The paths have no zeros, by \eqref{model:mesh-path-height}.  Equations
\eqref{model:mesh-path-height}--\eqref{model:mesh-path-length} imply
that the oscillation of \(C_0\) along any one of the three paths with
index \(j\) is at most
\begin{equation}
 e^{-2\lambda u_j+o(u_j)}.
 \label{model:mesh-oscillation}
\end{equation}
Choose an arbitrary point \(z_j\in\gamma_j\).  The connections just
proved join \(z_j\) to \(z_{j+1}\) by pieces of
\(\gamma_j,\sigma_j^-,\gamma_{j+1}\).  Thus
\(|C_0(z_{j+1})-C_0(z_j)|\le e^{-2\lambda u_j+o(u_j)}\).
The gaps \(u_{j+1}-u_j\) eventually exceed one.  For any fixed
\(\eta\in(0,2\lambda)\) and all sufficiently large \(j\), the
telescoping tail is therefore bounded by
\[
 \sum_{k\ge j}e^{-(2\lambda-\eta)u_k}
 \le \frac{e^{-(2\lambda-\eta)u_j}}
 {1-e^{-(2\lambda-\eta)}}.
\]
It follows that \(C_0(z_j)\) tends to a constant \(c_*\).  Letting
\(\eta\) tend to zero, and using
\eqref{model:mesh-oscillation} to reach every point of the three
paths, gives the uniform estimate
\begin{equation}
 |C_0-c_*|\le e^{-2\lambda u_j+o(u_j)}
 \quad\text{on }\gamma_j\cup\sigma_j^-\cup\sigma_j^+.
 \label{model:mesh-primitive-limit}
\end{equation}
This is one common constant for the whole mesh, not a separate
constant for each band or each side strip.

The function
\begin{equation}
 G=F(C_0-c_*)
 \label{model:mesh-vanishing-function}
\end{equation}
extends holomorphically through every zero of \(F\): its apparent
pole cancels, with limiting value \(-a^{-1}\) at the zero above.
The upper bound on the enlarged regions, together with
\eqref{model:mesh-primitive-limit}, gives
\begin{equation}
 |G|\le e^{-\lambda u_j+o(u_j)}\longrightarrow0
 \quad\text{uniformly on }
 \gamma_j\cup\sigma_j^-\cup\sigma_j^+.
 \label{model:mesh-G-bound}
\end{equation}

\paragraph{Enclosing every shallow point and reflecting.}
The arrangement of the paths is illustrated in \Cref{fig:mesh}.
\begin{figure}[t]
\centering
\begin{tikzpicture}[>=Stealth,x=1.55cm,y=1.2cm]
\draw[thick] (-3.1,3.05)--(3.1,3.05);
\node[font=\small,anchor=west] at (3.13,3.05) {$\Im z=0$};
\foreach \h in {.5,1.35,2.2}{
 \fill[deepblue!7] (-3,\h-.20) rectangle (3,\h+.20);
}
\draw[deepblue,thick] (-3,.46)--(-2.1,.51)--(-1.3,.43)--(-.5,.59)
 --(.2,.48)--(.9,.42)--(1.7,.55)--(2.2,.46)--(3,.52);
\draw[deepblue,thick] (-3,1.34)--(-2.2,1.42)--(-1.5,1.27)--(-.6,1.37)
 --(.2,1.29)--(1.1,1.44)--(1.9,1.32)--(3,1.37);
\draw[deepblue,thick] (-3,2.17)--(-2.1,2.28)--(-1.4,2.14)--(-.5,2.24)
 --(.3,2.13)--(1.0,2.26)--(1.8,2.15)--(3,2.23);
\draw[deepblue,thick] (-1.95,.29)--(-1.87,.82)--(-2.03,1.16)
 --(-1.92,1.71)--(-2.04,2.40);
\draw[deepblue,thick] (1.95,.29)--(2.03,.85)--(1.88,1.18)
 --(2.02,1.76)--(1.91,2.40);
\fill[warmred] (.45,1.75) circle (1.7pt);
\node[font=\small,anchor=west] at (.53,1.77) {target point};
\node[font=\small,anchor=east] at (-3.10,.50) {$u_j$};
\node[font=\small,anchor=east] at (-3.10,1.35) {$u_{j+1}$};
\node[font=\small,anchor=east] at (-3.10,2.20) {$u_{j+2}$};
\node[font=\small,align=center] at (0,-.15)
 {$\log|F|\ge\lambda u_j-o(u_j)$ on the connecting paths};
\end{tikzpicture}
\caption{The mesh in the pure first-field model. Vertical spacing is
schematic: the bands lie at logarithmic depths \(u_j+o(u_j)\).
Horizontal paths and side connectors enclose every sufficiently shallow
point of a smaller interval. Their high-modulus bounds give one common
limit \(c_*\) for the meromorphic primitive \(C_0\), where \(C_0'=F^{-2}\).
The product \(G=F(C_0-c_*)\) is holomorphic even at zeros of \(F\);
its boundary estimate therefore passes to the enclosed regions by the
maximum principle.}
\label{fig:mesh}
\end{figure}

Write \([\alpha_j,\beta_j]\) for a logarithmic-depth interval
containing the entire routed horizontal path \(\gamma_j\).  The
fixed-factor enlargements give
\(\alpha_j=u_j-s_j-O(1)\) and
\(\beta_j=u_j+s_j+O(1)\), and
\(\beta_j<\alpha_{j+1}\) for large \(j\).
Given \(z=x-\ii e^{-v}\) with \(x\in[-1,1]\) and \(v\) large,
choose \(j\) so that
\[
 \beta_j<v\le\beta_{j+1}.
\]
Then \(\gamma_j\) lies strictly below \(z\), whereas
\(\gamma_{j+2}\) lies strictly above \(z\).  The two paths
\(\sigma_j^-\) and \(\sigma_j^+\) connect across both bands and
lie strictly to the left and right of \(z\).

Choose an intersection of each side path with the lower and upper
horizontal paths, and use the corresponding path segments to form a
closed walk, in the order lower, right, upper, left.  More explicitly,
along either side path take its last lower-path hit before its first
subsequent upper-path hit; its intervening segment then joins those
two paths.  Join these endpoints along the horizontal paths.  Every
piece lies in the indicated open half-plane relative to \(z\), and
the four corner points lie in the corresponding quadrants.  This
walk has winding number one about \(z\): replace each of its pieces
by a straight segment within its open half-plane, which is convex
and avoids \(z\).  The resulting quadrilateral traverses successively
the four quadrants and has total change of argument \(2\pi\).
Equivalently, the arguments on the lower, right, upper, and left
pieces can be chosen in their respective half-plane intervals and
then joined in this order.  Self-intersections of a path or of the
walk do not affect this winding calculation.

The complementary component containing \(z\) is therefore bounded
and has boundary in this walk.  It lies compactly in the half-plane.
The maximum principle for the holomorphic function \(G\) gives its
bound there from the walk.  The upper horizontal piece has index
\(j+2\), and \(u_{j+2}=u_j+o(u_j)\), so
\eqref{model:mesh-G-bound} yields
\[
 |G(x-\ii e^{-v})|\le e^{-\lambda u_j+o(u_j)}.
\]
The same estimate is immediate if the point lies on a mesh path.
As \(v\to\infty\), the selected \(j\) tends to infinity uniformly
for \(x\in[-1,1]\).  Hence \(G\) has continuous boundary value zero
on \((-1,1)\), uniformly on compact subintervals.

Schwarz reflection across this interval extends \(G\) holomorphically
by \(\overline{G(\overline z)}\).  The extension vanishes on the
interval, so the identity theorem gives \(G\equiv0\).  On any disk
where \(F\ne0\), this says \(C_0\equiv c_*\), contradicting
\(C_0'=F^{-2}\ne0\).  The contradiction proves the Theorem.
\end{proof}

\begin{proof}[Proof of \Cref{thm:sequence-exclusion}]
Assume that a sequence in \Cref{def:prohibited-sequence} exists.
Its primary moments have a limit by
\Cref{model:moment-compactness}.  Apply
\Cref{lem:contact-dichotomy} to its logarithmic law.
The steep-contact alternative contradicts
\Cref{thm:exponential-obstruction}, as proved in
\Cref{model:steep-excluded}.  The bounded-contact
alternative gives the fields in \Cref{model:selected-fields}.
Their nonzero first field has holomorphic quotient \(F''/F\).
The limiting impact equation and parameter continuation,
\Cref{prop:pure-model,prop:model-growth}, give its uniform
upper growth.  \Cref{model:normal-epochs,prop:model-witnesses}
give common relatively dense epochs and uniform witnesses
at every horizontal center.  These are precisely the
hypotheses contradicted by \Cref{thm:mesh-contradiction}.
Both alternatives are impossible, proving the Theorem.
\end{proof}

\begin{corollary}[The critical angle budget is available]
\label{model:angle-budget-activated}
In the remaining finite-wall critical case, the angle budget
of \Cref{thm:angle-budget-reduction} and the polynomial bounds
of \Cref{cor:critical-polynomial-bounds} hold unconditionally.
In particular, for some \(\epsilon>0\),
\[
 \int_{k_0}^K\mu_f(k)^2\,\dd k\le C K^{2-\epsilon}.
\]
\end{corollary}

\begin{proof}
The sole forward assumption in the critical reduction was the
nonexistence of the sequences in
\Cref{def:prohibited-sequence}.  This is
\Cref{thm:sequence-exclusion}.  Apply the two stated
results.  Their conclusions were not used anywhere in the
proof of that theorem.
\end{proof}
